\chapter{先指定问题、模型和证据边界}
\section{这次扩大到哪里}
研究对象是非相对论磁重联中离子与电子的能量转化、温升和热输运。主要纳入有独立流体演化或可明确应用到流体方程的闭合：共速双温碰撞流体、双温Hall及惯性扩展、五矩、十矩、压力各向异性状态方程、Landau/陀螺流体，以及按来源分群的电子多流体。两种矩阶数与两种电磁约化可以交叉组合；“十矩”不自动意味着全Maxwell，“全Maxwell”也不自动意味着有正确热流。

离子--中性“两流体”、相对论电子--正电子两流体不作为本轮的离子--电子主线；并非它们不重要，而是物种身份与目标能量账本不同。PIC、Vlasov和混合模型只承担必要的参照或闭合观测检验，不扩展为另一份动理学综述。论文用双流体方程诊断PIC，也不等于独立推进了双流体模型。

检索使用网页检索、出版机构、arXiv、MIT与OSTI等公开来源，结合本库原文和引文追踪。关键词覆盖 two-fluid/ten-moment/twenty-moment/anisotropic pressure/heat flux/ion electron heating/reconnection，以及MAST、HIT-SI、MMS。逐次查询与访问状态保存在\path{search/}。搜索命中和题名筛查不是全文核读；本轮没有穷尽数据库、递归穷尽引文，也没有把未取得的正文计为新读。

旧库的全文记录按稳定键复用，定向复查不增加计数。新增读到的接收稿保留原文件、哈希及物理页码；正文、方法、讨论、图表及文内适用附录均读完，并直接查看关键PDF页后，才进入新核读记录。前序方法或独立补充尚未取得时明确保留缺口。报告中的“未报告”指本次所读版本没有提供足以支持该判断的信息，不意味着作者从未做过相应工作。

\section{Hall基线不能只有一个名字}
\begin{longtable}{P{16mm}P{54mm}P{78mm}}
\toprule 基线 & 必须写明的方程 & 哪类比较才有意义\\\midrule\endfirsthead
\toprule 基线 & 必须写明的方程 & 哪类比较才有意义\\\midrule\endhead
H0 & 等温或规定$p_e,p_i$的Hall & 温度不独立演化；任何温升预测先需要能量方程。\\
H1 & Hall欧姆律、总压或总能，固定物种分配 & 能比较总内能，但不能把指定的$T_i/T_e$当输出。\\
H2 & 独立标量$p_i,p_e$；各自对流；物理热源与输运 & 可比较温度、交换和损失；仍不能输出完整压力张量。Browning与HIT-SI是具体实现，不是统一标准。\\
H3 & H2加电子惯性和明确能量一致性 & 可比较电子动能储存与惯性破冻结；仍无自动产生的不可逆电子热源。增强电子质量也是参数。\\
H4 & Hall/惯性约化加各向异性或完整压力演化 & 与十矩的差别须再检查密度约束、位移电流、热流及碰撞，不能把压力自由度的成功归为“战胜所有Hall”。\\
\bottomrule\end{longtable}

这里H0--H4是本报告的比较标签，并非文献正式命名。H2和H3并非严格嵌套的现成代码：热源、壁面和电子惯性的实现可能同时变化。将两份不同几何的计算排在一张图上，只能构成跨条件启发，不能当成控制变量实验。

\section{成功证据分成四个层级}
第一层是模型能力：方程包含所需自由度，例如十矩能表示$P_{xy}$。第二层是数值可信度：给定闭合下，目标观测量随网格、步长和数值限制器收敛。第三层是模型比较：在匹配参数下，一个闭合比另一个更接近参照。第四层才是观测验证：经过测量响应、采样和不确定度处理后，与实验一致。层级之间没有自动推论。

对加热还须多一道判别：测量的温度峰、积分内能、焓流、电场沉积、模型熵产或分布函数尾分别是什么。相近重联率只说明磁通转移的一项指标相近；温度峰相近可能由不同热源和不同损失相抵得到。两者都不足以独自证明机制。

本报告将已解决的表述限定为“某一条件下某一可观测量已得到有证据的解释”。“模型表示不了”“闭合不可靠”“没有做同参数对照”“原文未取得”分别是不同缺口。读者应据此判断下一步是增添方程、改善闭合、设计对照，还是继续取得证据。
